\documentclass[10pt,a4paper]{amsart}
\usepackage[margin=19mm]{geometry}
\usepackage{amsmath,amssymb,mathtools}
\usepackage[scr=rsfs,cal=euler]{mathalfa}
\usepackage{xfrac}
\usepackage[unicode,hidelinks,bookmarks=false]{hyperref}
\newcommand{\ie}{i.e.\ }
\newcommand{\cf}{cf.\ }
\newcommand{\resp}{respectively\ }
\newcommand{\Subsection}{\subsection}
\providecommand{\nobreakdash}{\nobreak\hbox{-}}
\setlength{\emergencystretch}{2em}
\multlinegap0pt
\usepackage{mdframed}
\usepackage{mdframed}
\usepackage{mdframed}
\usepackage{mdframed}
\usepackage[shortlabels]{enumitem}
\usepackage{xcolor}
\usepackage{dsfont}
\usepackage{amssymb}
\usepackage{mdframed}
\usepackage{mathtools}
\usepackage{siunitx}
\usepackage{algorithm}\usepackage{algpseudocode}
\usepackage{breqn}
\newtheorem{assumption}{Assumption}
\newtheorem{lemma}{Lemma}
\newtheorem{theorem}{Theorem}
\title{Linesearch Method for \\ Accelerated Composite Minimization}
\author{Reza Rahimi Baghbadorani, Sergio Grammatico, Peyman Mohajerin Esfahani}
\date{Source: arXiv:2405.03414v3 (CC BY 4.0)}
\begin{document}
\maketitle

\noindent\emph{Source and licence.} Linesearch Method for \\ Accelerated Composite Minimization. Version arXiv:2405.03414v3 (CC BY 4.0), distributed by arXiv under the Creative Commons Attribution 4.0 International licence, http://creativecommons.org/licenses/by/4.0/, as recorded in the arXiv licence metadata for this exact version and shown on the abstract page https://arxiv.org/abs/2405.03414v3. Full licence notice: \texttt{licenses/arxiv-2405.03414-CC-BY-4.0.txt}.\par\smallskip

\noindent\textit{Editorial scope.} Counted unit: the article's theorem theorem1 (source0.tex lines 521--531). The article's complete original proof of it is delivered with it (source0.tex lines 532--627, one \texttt{proof} environment of 96 lines). No mathematics is written, repaired or re-derived: the statement and the proof are the article's own lines, in the article's own notation, and no retained line is substituted or re-flowed.  Retained uncounted as the source-local context the proof needs: lines 165--167; lines 225--227; lines 303--306; lines 313--316; lines 329--331; lines 334--361; lines 348--350.  Nothing was omitted from any retained span, so the package carries no omission record.  No retained line contains a citation, so the wrapper carries no bibliography and no bibliography entry.  No retained line contains a figure environment, an image-inclusion command or drawing code, so no image is delivered and the package declares no asset dependency.  Editorial formatting only, in addition: an AMS \texttt{amsart} preamble that restates the article's own declarations and macros used by the retained lines, namely the environments \texttt{assumption}, \texttt{lemma}, \texttt{theorem} on separate counters, together with the standard packages those lines need.  The retained environments print in this document as Assumption 1, Lemma 1, Theorem 1 in order of appearance, whereas the article prints the same items with the numbering of its own full document.  The complete native source of the article as distributed in the arXiv e-print for this exact version is bundled unchanged as \texttt{sources/158705/source0.tex}.
\begin{equation}\label{stepsize-search direction}
    x_{k+1} = x_k + \lambda_k d_k. %, \,\,\, k \in \mathbb{N} 
\end{equation}

            \begin{equation}\label{composite problem}
                \min_{x \in \mathcal{X}} F(x) = \min_{x \in \mathcal{X}} f(x)+h(x)
            \end{equation}

\begin{align}
    & \|\nabla f(x) - \nabla f(y)\| \leq L\|x-y\|, \label{eq2} \\
    & f(y) \leq f(x) + \langle \nabla f(x),y-x \rangle + \dfrac{L}{2} \|x-y\|^2. \label{eq1}
\end{align}

\begin{align}
\label{grad_mapping}
G^{f}_{\lambda h}(x) := \dfrac{1}{\lambda}\Big(x - \mathrm{prox}_{\lambda  h}\big(x - \lambda \nabla f(x)\big)\Big),   
\end{align}

\begin{assumption}[Convex regularity]\label{assum2}
The convex function~$F$ in \eqref{composite problem} has bounded level sets and a finite minimizer, where the convex terms~$f(x)$ and $h(x)$ are smooth and prox-friendly, respectively. 
\end{assumption}

\begin{lemma}[Gradient mapping]\label{keylemma}
Let $G^{f}_{\lambda h}(x)$ be the gradient mapping defined in \eqref{grad_mapping} for a smooth convex function $f$, a possibly nonsmooth function $h$, and a positive constant $\lambda$ in $\mathbb{R}_+$.
\begin{enumerate}[label=(\roman*), itemsep = 0mm, topsep = 0mm, leftmargin = 7mm]
\item \label{lema3}
{\bf Stationary condition:} The vector $x$ is a minimizer of \eqref{composite problem} if and only if $G^{f}_{\lambda h}(x) = 0$.

\item \label{lema1}
{\bf Convexity-like inequality:} For all $x,y$ in $\mathbb{R}^d$ we have
\begin{align*}
    h\left(x-\lambda G^{f}_{\lambda h}(x)\right) \leq h(y) - \langle G^{f}_{\lambda h}(x) - \nabla f(x), y - (x - \lambda G^{f}_{\lambda h}(x)) \rangle.  %\quad \forall x,y \in \mathbb{R}^d
\end{align*}

\item \label{lema2}
{\bf Increment bound:} Considering $x^+ = x - \lambda G^{f}_{\lambda h}(x)$, for any $z \in \mathbb{R}^d$ we have
\begin{equation*}\label{eq3}
F(x^+) - F(z) \leq \langle x^+ - x,  \nabla f(x^+) - \nabla f(x) + \dfrac{1}{2}G^{f}_{\lambda h}(x) \rangle - \Big( \dfrac{1}{2\lambda} \|x^+ - x\|^2 + \dfrac{1}{\lambda}\langle x^+ -x, x - z \rangle \Big).
\end{equation*}

\item \label{lema4linesearch}
{\bf Zero-order linesearch:} Let the update direction~$d_k = -G^{f}_{\lambda h}(x)$. If $\lambda$ satisfies
%If $\lambda$ satisfies the following inequality with search direction $d_k = -G^{f}_{\lambda h}(x)$
%To provide the convergence analysis, we need the first term on the right-hand side of \ref{lema2} to remain negative. The following stepsize rule guarantees the negativity of this term.
\begin{equation}\label{ZO linesearch}
    \phi_k(2\lambda) \leq \phi_k(\lambda) - \lambda \langle G^{f}_{\lambda h}(x), \nabla f(x) \rangle + \dfrac{\lambda}{2}\|G^{f}_{\lambda h}(x)\|^2,
\end{equation}
we then have $\langle x^+ - x,  \nabla f(x^+) - \nabla f(x) + \dfrac{1}{2}G^{f}_{\lambda h}(x) \rangle \leq 0$.
\end{enumerate}
\end{lemma}

\begin{equation*}
F(x^+) - F(z) \leq \langle x^+ - x,  \nabla f(x^+) - \nabla f(x) + \dfrac{1}{2}G^{f}_{\lambda h}(x) \rangle - \Big( \dfrac{1}{2\lambda} \|x^+ - x\|^2 + \dfrac{1}{\lambda}\langle x^+ -x, x - z \rangle \Big).
\end{equation*}

\begin{theorem}[Non-accelerated adaptive stepsize]\label{theorem1}
Consider the function~$F$ in \eqref{composite problem} as a composition of a smooth function $f(x)$ and a prox-friendly function $h(x)$. Assume that the sequence $\left (x_k \right)_{k \in \mathbb{N}}$ generated by the update algorithm \eqref{stepsize-search direction} in which the direction and stepsize rule are
\begin{equation*}\label{alg:Algorithm1}
\tag{Alg.\,1}
\begin{cases}
  \lambda_k = \max\limits \left \{\lambda \in \mathbb{R}_+ \,|\,\, \phi_k(2\lambda) \leq \phi_k(\lambda) - \lambda \langle G^{f}_{\lambda h}(x_k), \nabla f(x_k) \rangle + \dfrac{\lambda}{2}\|G^{f}_{\lambda h}(x_k)\|^2\right\},\\      
  d_k = -G^{f}_{\lambda_k h}(x_k).\\
\end{cases}
\end{equation*}
Then, we have the uniform stepsize lower bound $\lambda_k \geq {1}/{3L}$ where $L$ is the smoothness constant of $f$. Moreover, we have the sub-optimality bound $F(x_k) - F(x^*) \leq \dfrac{D}{k}$ where $D$ is a constant that only depends on the initial condition $x_0$, the optimal solution $x^*$, and the smoothness constant $L$.
\end{theorem}

\begin{proof}
First, to show that $\lambda_k$ is bounded from below by ${1}/{3L}$, we just need to show that the condition in the stepsize rule always is satisfied by $\lambda_k = 1/{3L}$. Considering the smoothness of $f$, we can write the inequality \eqref{eq1} for the pair $(x,x^+)$ where $x^+ = x - \lambda G^{f}_{\lambda h}(x)$, which arrives at the two inequalities
\begin{align*}
f(x^+) &\leq f(x) - \langle \nabla f(x),\lambda G^{f}_{\lambda h}(x) \rangle + \dfrac{L\lambda ^2}{2} \|G^{f}_{\lambda h}(x)\|^2, \\
f(x) &\leq f(x^+) + \langle \nabla f(x^+),\lambda G^{f}_{\lambda h}(x) \rangle + \dfrac{L\lambda ^2}{2} \|G^{f}_{\lambda h}(x)\|^2.
\end{align*}
Adding the two sides of the above inequalities leads to
\begin{equation}\label{eq4}
\langle G^{f}_{\lambda h}(x), \nabla f(x^+) - \nabla f(x) + {L\lambda}G^{f}_{\lambda h}(x) \rangle \geq 0, \qquad \forall \, \lambda \in \mathbb{R}_+, \, x \in \mathbb{R}^n.   
\end{equation}
%%%%%%%%%%%%%
{By the smoothness of the scalar function $\phi_k(\lambda) = f(x_k + \lambda d_k)$, it follows that
\begin{align*}
    f(x_k + 2\lambda d_k)
    \leq f(x_k + \lambda d_k)
    + \nabla f(x_k + \lambda d_k)^\top \bigl( (x_k + 2\lambda d_k) - (x_k + \lambda d_k) \bigr)
    + \frac{L}{2}\|\lambda d_k\|^2.
\end{align*}
Rewriting the above inequality yields
\begin{align*}
    f(x_k + 2\lambda d_k) - f(x_k + \lambda d_k)
    \leq \lambda\, d_k^\top \nabla f(x_k + \lambda d_k)
    + \frac{L\lambda^2}{2}\|d_k\|^2.
\end{align*}
To ensure that the stepsize rule in \ref{alg:Algorithm1} is satisfied, and using the definition of the direction $d_k$, we require
{\color{black}
\begin{align*}
    - \lambda G^{f}_{\lambda h}(x_k)^\top
    \nabla f\bigl(x_k - \lambda G^{f}_{\lambda h}(x_k)\bigr)
    + \frac{L\lambda^2}{2}\|G^{f}_{\lambda h}(x_k)\|^2
    \leq
    - \lambda \langle G^{f}_{\lambda h}(x_k), \nabla f(x_k) \rangle
    + \frac{\lambda}{2}\|G^{f}_{\lambda h}(x_k)\|^2.
\end{align*}
Rearranging the terms and dividing both sides by $\lambda>0$, this inequality is equivalent to}
\begin{align}\label{thr2_3_ineq1}
    0 \leq
    \left\langle G^{f}_{\lambda h}(x_k),
    \nabla f\bigl(x_k - \lambda G^{f}_{\lambda h}(x_k)\bigr) - \nabla f(x_k)
    \right\rangle
    + \frac{1 - \lambda L}{2}\|G^{f}_{\lambda h}(x_k)\|^2.
\end{align}
Comparing this expression with inequality~\eqref{eq4}, which holds universally, we deduce that the inequality \eqref{thr2_3_ineq1} holds whenever
\begin{align}\label{thr2_3_ineq2}
    \lambda L \leq \frac{1 - \lambda L}{2},
\end{align}
which implies $\lambda \leq \frac{1}{3L}$.}
%%%%%%%%%%%%%%
% Due to the smoothness of $\phi(\cdot)$, the stepsize rule in \ref{alg:Algorithm1} is satisfied whenever
% \begin{align*}
%     \big\langle G^{f}_{\lambda h}(x),\, \nabla f(x^+) - \nabla f(x) + \tfrac{1 - L\lambda}{2} \, G^{f}_{\lambda h}(x) \big\rangle \geq 0,
% \end{align*}
% which coincides with the inequality condition in \eqref{eq4} when $\lambda = 1/3L$, showing that the stepsize rule in \ref{alg:Algorithm1} is always satisfied by choosing $\lambda = 1/3L$.
Next, we prove the convergence of \ref{alg:Algorithm1}. By putting $z = x_k$ in the inequality \ref{eq3} we have
\begin{equation}\label{conv-analysis}
F(x_{k+1}) - F(x_k) \leq - \dfrac{\lambda_k}{2} \|G^{f}_{\lambda_k h}(x_k)\|^2  
\end{equation}
\textcolor{black}{Summing the previous inequality from $k=0$ to $T-1$ and utilizing the stepsize lower bound $\lambda_k \geq 1/(3L)$, we obtain
\begin{equation*}
\frac{1}{6L}\sum_{k=0}^{T-1} \|G^{f}_{\lambda_k h}(x_k)\|^2 \leq F(x_0) - F(x_{T}).
\end{equation*}
By Assumption \ref{assum2}, the objective function $F(x)$ is bounded from below by its optimal value $F^*$. Substituting this lower bound and taking the limit as $T \to \infty$ yields
\begin{equation*}
\sum_{k=0}^\infty \|G^{f}_{\lambda_k h}(x_k)\|^2 \leq 6L (F(x_0) - F^*) < \infty.
\end{equation*}
Because the sequence of function values is monotonically decreasing and bounded from below, the non-negative infinite series on the left-hand side converges to a finite value. This implies that the elements of the series must asymptotically approach zero. Thus, we conclude that
\begin{equation*}
\lim_{k \to \infty} \|G^{f}_{\lambda_k h}(x_k)\|^2 = 0,
\end{equation*}
which satisfies the standard first-order optimality conditions and establishes convergence by Lemma~\ref{keylemma}~\ref{lema3}.}
% Telescoping the above inequalities from $k=0$ to $k=\infty$ and using the fact that $\lambda_k \geq 1/3L$ gives us
% \begin{equation*}
%     F(x_\infty) - F(x_0) \leq -\dfrac{1}{6L}\sum_{i=0}^\infty \|G^{f}_{\lambda_i h}(x_i)\|^2
% \end{equation*}
% According to Assumption \ref{assum2} that $F(x)$ has a finite minimizer, we conclude $\|G^{f}_{\lambda h}(x_\infty)\|^2 \to 0$ which shows the convergence from \ref{lema3}. 
To derive the rate of convergence we start from \ref{lema2} with $z = x^*$.
\begin{equation}\label{aux-conv}
   F(x_{k+1}) \leq F(x^*) - \dfrac{1}{2\lambda_k} \|x_{k+1} - x_k\|^2 - \dfrac{1}{\lambda_k}\langle x_{k+1} - x_k, x_k -  x^* \rangle  
\end{equation}
By adding and subtracting $\dfrac{1}{\lambda_k} \langle x_{k+1}, x_k \rangle$ and $\dfrac{1}{\lambda_k} \langle x_k, x_k \rangle$ terms and expanding the right-hand side of \eqref{aux-conv} we have  
\begin{equation}\label{aux-conv2}
     F(x_{k+1}) \leq F(x^*) + \dfrac{1}{2\lambda_k}\Big(\|x_k - x^*\|^2 - \|x_{k+1} - x^*\|^2\Big)
\end{equation}
{\color{black} Using \eqref{aux-conv2} and the lower bound of $\lambda_k$, we can write
\begin{equation*}
    F(x_{k+1}) \leq F(x^*) + \frac{3}{2}L\Big(\|x_k - x^*\|^2 - \|x_{k+1} - x^*\|^2\Big)
\end{equation*}
Summing the above inequality from $k = 0$ to $k = T-1$ gives us
\begin{equation*}
    \sum_{k=1}^{T}\Big(F(x_k) - F(x^*)\Big) \leq \frac{3}{2}L\Big(\|x_0 - x^*\|^2 - \|x_T - x^*\|^2\Big) \leq \frac{3}{2}L\|x_0 - x^*\|^2
\end{equation*}
Because the sequence of $F(x_k)$ is nonincreasing, averaging yields
\begin{align*}
    F(x_T) - F(x^*) \leq \dfrac{1}{T}\sum_{k=1}^{T}\Big(F(x_k) - F(x^*)\Big) \leq \frac{3}{2T}L\Big(\|x_0 - x^*\|^2 - \|x_T - x^*\|^2\Big) \leq \dfrac{3L\|x_0 - x^*\|^2}{2T} = \dfrac{D}{T}.
\end{align*}}
\end{proof}

\end{document}
